\chapter*{Реферат}
\thispagestyle{plain}

Пояснительная записка содержит 33 страницы, 4 главы, 3 приложения, 17 использованных источников.

Ключевые слова: Retrieval-Augmented Generation, RAG, большие языковые модели, векторный поиск, Qdrant, embeddings, корпоративная база знаний, вопросно-ответная система, dense retrieval.

Целью работы является исследование и разработка базового прототипа RAG-ассистента, обеспечивающего поиск релевантных фрагментов в корпусе русскоязычных бизнесово-технических документов и генерацию ответа на пользовательский вопрос с использованием найденного контекста.

В рамках работы рассматривается первая MVP-итерация системы. Она сознательно ограничена базовой схемой Retrieval-Augmented Generation: пользовательский вопрос кодируется embedding-моделью, по полученному вектору выполняется поиск ближайших фрагментов в Qdrant, найденные чанки объединяются в контекст, после чего локальная языковая модель формирует ответ на русском языке и возвращает список источников. Такой вариант позволяет проверить основной путь от исходных документов до ответа без включения дополнительных механизмов ранжирования и агентной логики.

В первой главе выполнен анализ предметной области и проблемы поиска информации в корпоративной базе знаний. Рассмотрены ручной поиск, полнотекстовый поиск, семантический поиск на основе embeddings и RAG-подход. Показано, что для корпуса корпоративных документов требуется не только нахождение файла, но и получение прикладного ответа с указанием источников.

Во второй главе рассмотрены теоретические основы больших языковых моделей, векторного представления текста, векторных баз данных и базовой архитектуры RAG. Отдельно выделены offline-этап подготовки индекса и online-этап обработки пользовательского вопроса.

В третьей главе спроектирован базовый RAG-ассистент для корпоративной базы знаний. Описаны функциональные и нефункциональные требования, архитектура обработки документов, организация dense retrieval в Qdrant и схема генерации ответа по найденному контексту.

В четвертой главе приведено описание реализации MVP и первичной проверки работоспособности. По отчетам проекта обработано 36 документов форматов DOCX, PDF и XLSX; после парсинга получено 2300 структурированных блоков общим объемом 589074 символов; после чанкинга сформирован 481 фрагмент. Для построения embeddings используется модель \texttt{deepvk/USER-bge-m3} размерности 1024, для векторного хранилища применяется Qdrant, для генерации ответа используется локальная LLM \texttt{qwen2.5:7b-instruct} через OpenAI-compatible API. По отчету \path{data/reports/retrieval/retrieval_eval_report.json} первичная проверка базового dense retrieval на 150 вопросах показала \texttt{Hit@5=0.7467}, \texttt{MRR@5=0.5676}, \texttt{Recall@5=0.6522}.

%%% Local Variables:
%%% TeX-engine: xetex
%%% End:
